% ============================================================================
%  NotFlex: Compilador de Expresiones Regulares a Código C Optimizado
%  para Entornos Embebidos
%  Formato: Springer LNCS (clase llncs) -- Ingeniería en Informática, UGD
%  Compilar con pdfLaTeX (Overleaf incluye llncs.cls por defecto).
%  VERSIÓN CON TEXTOS REFORMULADOS (tono simple, de alumno)
% ============================================================================
\documentclass[runningheads]{llncs}

\usepackage[T1]{fontenc}
\usepackage[utf8]{inputenc}
\usepackage[spanish,es-tabla,es-nodecimaldot,es-noshorthands]{babel}
\usepackage{amsmath,amssymb}
\usepackage{graphicx}
\usepackage{booktabs}
\usepackage{tabularx}
\usepackage{array}
\usepackage{xcolor}
\usepackage{listings}
\usepackage{tikz}
\usetikzlibrary{arrows.meta,positioning,fit}
\usepackage[hidelinks]{hyperref}

% --- Etiquetas en castellano compatibles con llncs ---------------------------
\renewcommand\keywordname{{\bfseries Palabras clave:}}
\renewcommand\andname{y}               % "A y B"
\renewcommand\lastandname{\unskip\ y}  % "A, B y C"
\renewcommand{\lstlistingname}{Listado}

% --- Estilo de código C ------------------------------------------------------
\lstdefinestyle{cstyle}{
  language=C,
  basicstyle=\ttfamily\scriptsize,
  keywordstyle=\bfseries\color{blue!55!black},
  commentstyle=\itshape\color{black!55},
  stringstyle=\color{red!50!black},
  numbers=left,
  numberstyle=\tiny\color{black!45},
  numbersep=6pt,
  frame=tb,
  framerule=0.4pt,
  xleftmargin=1.6em,
  framexleftmargin=1.6em,
  aboveskip=0.8\baselineskip,
  belowskip=0.6\baselineskip,
  columns=fullflexible,
  keepspaces=true,
  showstringspaces=false,
  upquote=true,
  literate={á}{{\'a}}1 {é}{{\'e}}1 {í}{{\'i}}1 {ó}{{\'o}}1 {ú}{{\'u}}1 {ñ}{{\~n}}1
}

\begin{document}

% ============================================================================
%  Frontmatter
% ============================================================================
\title{NotFlex: Compilador de Expresiones Regulares a Código C Optimizado para Entornos Embebidos}
\titlerunning{NotFlex}

\author{Hillebrand Giuliano \and Sienra Daiana \and González Martín}


\institute{Universidad Gastón Dachary, Posadas, Misiones, Argentina}

\maketitle

\begin{abstract}
NotFlex es  un compilador que toma una expresión regular y genera un autómata finito determinista y código C (estándar C99) listo para ser utilzado en sistemas embebidos e IoT. El código C generado es 2,2 y 11,5 veces más rápido que usar \texttt{regex.h} de POSIX dado que no utiliza memoria dinámica ni librerías externas. Además, el código C generado puede recorrer una cadena en tiempo lineal $\mathcal{O}(n)$. El compilador se divide en frontend y backend. En el frontend donde se realiza el análisis léxico, sintáctico, semántico y donde se arma el Árbol de Sintaxis Exacta; y en el backend el motor de autómatas aplica la construcción de Thompson, subconjuntos que agrupan caracteres en clases de equivalencia y minimización de Hopcroft.

\keywords{Compiladores \and Expresiones Regulares \and Autómatas Finitos \and Sistemas Embebidos \and Minimización de Hopcroft \and Código C Optimizado.}
\end{abstract}

% ============================================================================
\section{Introducción}
% ============================================================================
El foco de nuestro proyecto son los sistemas embebidos e IoT (Internet de las Cosas). En este contexto, la memoria y la capacidad de procesamiento son recursos limitados, y en ocasiones se necesita validar cadenas de texto, como tramas obtenidas por sensores, identificadores o comando por puerto serie. Para abordar estos problemas existen dos maneras de hacerlo, pero ambas presentan sus desventajas:
\begin{itemize}
    \item Escribir un validador a mano, a base de \texttt{if/switch} anidados que pueden ser complejos y donde es fácil equivocarse.
    \item Implementar una librería de expresiones regulares, como \texttt{regex.h} de POSIX que ocupan mucho de estos escasos recursos recién mencionados, porque incluso utilizan asignación de memoria dinámica. Además, si la librería implementada utiliza backtracking, como PCRE2, podría llegar a tardar un tiempo exponencial al utilizarse con ciertos patrones.
\end{itemize}

Como solución ante estas dificultades, proponemos NotFlex. Un compilador que hace el trabajo pesado en tiempo de compilación y dada una expresión regular de entrada, genera como salida una función C en el estándar C99 muy ligera que recorre una tabla de transiciones estática, por lo que no utiliza asignación dinámica de memoria y permite dedicar los limitados recursos a otras necesidades, ya que no depende de librerías adicionales ni de un motor de expresiones regulares.

\section{Metodología de Desarrollo}\label{sec:metodologia}

NotFlex se desarrolló de forma asistida por inteligencia artificial (IA). Implementamos un enfoque de Desarrollo Dirigido por Especificaciones y una suite de tests que nos permitió auditar el funcionamiento del código generado por IA para luego ser integrado oficialmente al proyecto. Esto NO quiere decir que el trabajo es una caja negra, realizado en su totalidad por IA a través de un "Claude, hacelo, y no cometas errores". Los integrantes del equipo estuvimos muy involucrados en el desarrollo del proyecto y vale destacar que este informe es 100\% de nuestra autoría y 0\% hehco con IA. 
Podemos resumir nuestro flujo de trabajo en \emph{el equipo investiga y especifica, luego la IA implementa, el equipo testea y audita para y finalmente el equipo integra el desarrollo al proyecto entregable}.

\subsection{De la investigación, al informe y luego a la especificación}

A medida que se desarrollaba el cursado de la cátedra de Compiladores, el equipo continuamente detallaba este informe con las necesidades y limitaciones que nuestro trabajo integrador debía cumplir. Con lo visto en la cátedra (las fases
clásicas de un compilador) definimos la idea, el objetivo, el
alcance, las transformaciones del motor y los escenarios de prueba que
demuestran que el objetivo se cumple. Y a partir de este documento derivamos tareas acotadas y precisas (por ejemplo, el analizador léxico, las definiciones con
tabla de símbolos, el motor de autómatas), acompañadas de sus requisitos y
su criterio de aceptación que luego entregamos a la IA como especificación.

Las IAs utilizadas fueron:
\begin{itemize}
    \item Gemini 4 Argón: modelo de inteligencia artificial de frontera desarrollado por Google DeepMind, diseñado específicamente para ofrecer un razonamiento profundo dinámico, inteligencia de alto nivel y el máximo rendimiento en flujos de trabajo complejos y exigentes.
    \item Claude Opus 5.5: modelo insignia de Anthropic diseñado específicamente para el trabajo de conocimiento de alta complejidad y el desarrollo de software agéntico de larga duración.
\end{itemize}

\subsection{Los tests como mecanismo de auditoría}

Para establecer una tarea dentro del proyecto, el criterio de aceptación de cada una de ellas se expresó como pruebas automáticas implementadas con \texttt{pytest}. El equipo fue definiendo los tests iterativamente\footnote{Algunos tests fueron generados con IA debido a su naturaleza compleja o extensa, como el test que se encarga de hacerle fuzzing al simulador del AFD con 1500 cadenas aleatorias para poder validar su funcionamiento.} a medida que el desarrollo del proyecto progresaba, al final del desarrollo, La suite completa cuenta con 268 pruebas y combina cuatro tipos:
\begin{itemize}
  \item \textbf{Unitarias y de integración:} Son unitarios y de integracón porque verifican que cada pieza o parte haga lo que se supone que debe hacer, por si misma (unitaria) o conectada a otras (integración). Detectan errores de lógica local como un mensaje de error en una posición equivocada pero están limitados a casos que el equipo pensó de antemano y no pueden testear lo que a nadie se le ocurrió. 
  \item \textbf{De caracterización:} Fueron escritos sobre el código existente para establecer lo que el código hacía en un momento determinado y nos permitieron modificar el código existente sin accidentalmente romper algo. Su objetivo es detectar regresiones, por ejemplo, si en algún momento alguien cambiaba accidentalmente como se tokenizaba el caracter \\ el test fallaba. También están limitados a que si accidentalmente se definía un test y se configuraba erroneamente, el test podría dar como válido un error.
  \item \textbf{Diferenciales:} Nos permiten comparar la salida del compilador con la de una implementación independiente. Por ejemplo, hay un test que se encarga de traducir el AST a una expresión de la librería re y verifica 12 ejemplos predefinidos con cadenas aceptadas y rechazadas
  \item \textbf{De propiedades no funcionales:} Este tipo de test no verificaba el correcto funcionamiento del código, si no que nos permitieron validar nuestro objetivo de generar código C ligero y muy optimizado para entornos con recursos limitados. Este tipo de tests nos permitió realizar las afirmaciones en la Sección ~\ref{sec:pruebas}
\end{itemize}

% ============================================================================
\section{Arquitectura General del Compilador}
% ============================================================================
El sistema sigue una serie de pasos secuenciales y conectados, clásico de un compilador, dividido en tres etapas (Fig.~\ref{fig:pipeline}):

\begin{figure}[htbp]
\centering
\begin{tikzpicture}[
  node distance=4mm and 5mm,
  paso/.style={draw, rounded corners=2pt, align=center, font=\scriptsize,
               minimum height=8mm, text width=2.1cm},
  flecha/.style={-{Stealth[length=2mm]}, thick},
  etapa/.style={draw=black!45, dashed, rounded corners=4pt,
                inner xsep=5pt, inner ysep=10pt},
  titulo/.style={font=\scriptsize\bfseries, inner sep=1pt}]
  % Etapa de análisis (frontend)
  \node[paso] (in) {Archivo\\\texttt{.regex}};
  \node[paso, right=of in]  (lex) {Análisis léxico\\(tokens)};
  \node[paso, right=of lex] (sin) {Análisis sintáctico\\(AST)};
  \node[paso, right=of sin] (sem) {Análisis semántico\\(AST expandido)};
  \node[paso, dashed, above=3mm of sem] (tab) {Tabla de símbolos};
  \draw[flecha] (in) -- (lex);
  \draw[flecha] (lex) -- (sin);
  \draw[flecha] (sin) -- (sem);
  \draw[{Stealth[length=2mm]}-{Stealth[length=2mm]}, dashed] (tab) -- (sem);
  % Etapa de síntesis (motor de autómatas)
  \node[paso, below=14mm of lex] (tho) {Thompson\\(AFND)};
  \node[paso, right=of tho, text width=2.8cm] (sub) {Subconjuntos y clases\\de equivalencia (AFD)};
  \node[paso, right=of sub] (hop) {Hopcroft\\(AFD mínimo)};
  \draw[flecha] (sem.south) -- ++(0,-5mm) -| (tho.north);
  \draw[flecha] (tho) -- (sub);
  \draw[flecha] (sub) -- (hop);
  % Etapa de generación (backend)
  \node[paso, below=14mm of hop] (gen) {Código C99\\(\texttt{.c}/\texttt{.h})};
  \draw[flecha] (hop) -- (gen);
  % Marcos de cada etapa
  \node[etapa, fit=(in)(sem)(tab)] (e1) {};
  \node[etapa, fit=(tho)(hop)] (e2) {};
  \node[etapa, fit=(gen)] (e3) {};
  \node[titulo, anchor=north west] at (e1.north west) {Análisis};
  \node[titulo, anchor=north east] at (e2.north east) {Síntesis};
  \node[titulo, anchor=north east] at (e3.north east) {Generación};
\end{tikzpicture}
\caption{Pasos secuenciales de compilación de NotFlex en tres etapas.}
\label{fig:pipeline}
\end{figure}

\begin{enumerate}
  \item \textbf{Frontend (Análisis):} El analizador léxico separa la expresión en tokens, el parser descendente recursivo arma el AST respetando la precedencia de los operadores, y el análisis semántico resuelve las definiciones con nombre usando una tabla de símbolos. Cada error se informa con su línea y su posición.
  \item \textbf{Motor de Autómatas (Síntesis):} Pasa el AST a un AFND con Thompson y lo convierte en AFD con subconjuntos, agrupando en el mismo paso los 256 valores posibles de un byte en clases de equivalencia. Por último aplica Hopcroft, un método que reduce al mínimo el número de estados del AFD.
  \item \textbf{Backend (Generación de Código):} Genera el archivo C (\texttt{.c}/\texttt{.h}) con la tabla de transiciones y el bucle de validación.
\end{enumerate}

% ============================================================================
\section{Especificación Sintáctica y Léxica}
% ============================================================================
El compilador acepta expresiones regulares, es decir, lenguajes de Tipo~3 en la jerarquía de Chomsky. La sintaxis soportada es:
\begin{itemize}
  \item \textbf{Literales y Concatenación:} \texttt{ab}
  \item \textbf{Alternancia (Unión):} \texttt{a|b}
  \item \textbf{Clausuras:} Kleene (\texttt{a*}), Positiva (\texttt{a+}) y Opción (\texttt{a?})
  \item \textbf{Agrupamiento y Clases:} \texttt{(a|b)*}, \texttt{[a-zA-Z0-9\_]}
  \item \textbf{Comodín:} \texttt{.} coincide con cualquier carácter salvo el salto de línea.
  \item \textbf{Escapes de caracteres:} \texttt{\textbackslash n}, \texttt{\textbackslash t} y una barra seguida de un símbolo, que lo vuelve literal (\texttt{\textbackslash +}, \texttt{\textbackslash .}).
  \item \textbf{Referencias:} \texttt{\{NOMBRE\}} usa una definición con nombre (ver Sección~\ref{sec:semantica}).
\end{itemize}

\subsection{Gramática}
El parser descendente recursivo implementa la siguiente gramática, con una función por no terminal:
\begin{center}
\small
\begin{tabular}{@{}r@{\ $\rightarrow$\ }l@{}}
\textit{expr}    & \textit{subexpr} (\texttt{|} \textit{subexpr})$^*$ \\
\textit{subexpr} & \textit{factor} \textit{factor}$^*$ \\
\textit{factor}  & \textit{átomo} (\texttt{*} $\mid$ \texttt{+} $\mid$ \texttt{?})$^*$ \\
\textit{átomo}   & \texttt{CHAR} $\mid$ \texttt{CLASS} $\mid$ \texttt{REF} $\mid$ \texttt{(} \textit{expr} \texttt{)}
\end{tabular}
\end{center}
La precedencia, de mayor a menor, es: cuantificadores, concatenación y unión. El conjunto PRIMERO de todos los no terminales es $\{\texttt{CHAR}, \texttt{CLASS}, \texttt{REF}, \texttt{(}\}$, y en cada repetición los tokens que la continúan son disjuntos de su conjunto SIGUIENTE, por lo que la gramática es LL(1). El analizador léxico agrega un token de fin de entrada que cumple el rol del símbolo~\$.

\subsection{Análisis semántico y tabla de símbolos}\label{sec:semantica}
Para dar nombre a las subexpresiones, la entrada admite definiciones al estilo de Flex, separadas del patrón principal por \texttt{\%\%}:
\begin{lstlisting}[basicstyle=\ttfamily\scriptsize, frame=tb, columns=fullflexible]
DIGITO      = [0-9]
ENTERO      = {DIGITO}+
NUMERO      = {ENTERO}(\.{ENTERO})?
TEMPERATURA = TEMP={NUMERO}
HUMEDAD     = HUM={ENTERO}
%%
{TEMPERATURA};{HUMEDAD}
\end{lstlisting}
La fase semántica registra cada definición en la tabla de símbolos y rechaza las definiciones duplicadas, las referencias sin definición, las referencias circulares (informando el camino, por ejemplo \texttt{A -> B -> A}) y los caracteres fuera del alfabeto de 256. Las definiciones que no se usan generan una advertencia. Su resultado es el AST del patrón principal con las referencias ya expandidas, que es la entrada del motor de autómatas.


% ============================================================================
\section{Limitaciones y Trabajo Futuro}
% ============================================================================
\begin{itemize}
  \item \textbf{Limitaciones:} Solo sirve para patrones que se conocen al compilar. Tampoco soporta negación de clases (\texttt{[\textasciicircum{}a-z]}), cuantificadores acotados (\texttt{a\{3,5\}}), anclas (\texttt{\textasciicircum}, \texttt{\$}), clases abreviadas (\texttt{\textbackslash d}, \texttt{\textbackslash w}) ni retroreferencias (estas últimas no se pueden hacer con lenguajes regulares). Todas ellas producen un error léxico en lugar de interpretarse como literales, para que el comportamiento coincida con el del módulo \texttt{re} de Python.
  \item \textbf{Trabajo Futuro:} (i)~Generar el código con \texttt{switch/case} como alternativa a las tablas; (ii)~Juntar varios patrones en un solo autómata, para llegar a un generador de analizadores léxicos completo.
\end{itemize}

% ============================================================================
\section{Transformaciones e Implementación del Motor}
% ============================================================================
El motor aplica tres transformaciones, una después de la otra (junto con la construcción del AST y la generación de C, el pipeline completo tiene cinco):
\begin{enumerate}
  \item \textbf{Construcción de Thompson:} Recorre el AST de forma recursiva y arma el AFND, uniendo los pedazos con transiciones épsilon ($\epsilon$).
  \item \textbf{Construcción de Subconjuntos con Clases de Equivalencia:} Antes de construir el AFD, junta los caracteres que se comportan igual en todo el autómata. Para un identificador de C (\texttt{[a-zA-Z\_][a-zA-Z0-9\_]*}) los 256 valores de un byte quedan en solo 3 clases (\emph{letra}, \emph{dígito}, \emph{otro}). Después saca el no determinismo calculando las $\epsilon$-clausuras, pero recorriendo esas pocas clases en lugar de 256 caracteres, y arma una tabla de transiciones más chica.
  \item \textbf{Minimización de Hopcroft:} Junta los estados que no se pueden distinguir. Cada estado que se elimina achica el arreglo que se genera.
\end{enumerate}

% ============================================================================
\section{Generación de Código C (Backend)}\label{sec:backend}
% ============================================================================
El backend genera en código C estándar (C99) la tabla de transiciones y el recorrido de la cadena. Un ejemplo, para un identificador, podría verse así:

\begin{lstlisting}[style=cstyle]
/* Mapeo de los 256 valores de un byte a 3 clases de equivalencia */
static const unsigned char clase[256] = { /* generado automáticamente */ };

/* Tabla de transiciones: [estado][clase] */
static const unsigned char trans[3][3] = {
    /*          letra  digito  otro */
    /* q0 */ {    1,     2,     2 },  /* q0: estado inicial */
    /* q1 */ {    1,     1,     2 },  /* q1: estado de aceptación */
    /* q2 */ {    2,     2,     2 }   /* q2: estado de rechazo */
};

int validar_id(const char *s) {
    int estado = 0;
    while (*s != '\0') {
        estado = trans[estado][clase[(unsigned char)*s++]];
    }
    return estado == 1;
}
\end{lstlisting}

La función no usa \texttt{malloc} (asignación de memoria), no depende de nada externo y hace un solo acceso a la matriz por cada carácter.

% ============================================================================
\section{Escenarios de Prueba y Evaluación} \label{sec:pruebas}
% ============================================================================
Para comprobar que el compilador funciona bien armamos pruebas que se pueden repetir. Todas forman parte de la suite de \texttt{pytest} (264 pruebas, todas aprobadas) y las dos últimas también se ejecutan con \texttt{make valgrind} y \texttt{make bench}:
\begin{enumerate}
  \item \textbf{Pruebas Unitarias:} Cubren el análisis léxico, sintáctico y semántico, los mensajes de error con su posición, la interfaz de línea de comandos y cada transformación del motor (por ejemplo, que el identificador de C produzca exactamente las 3 clases y la tabla de $3 \times 3$ de la Sección~\ref{sec:backend}). Los doce ejemplos incluidos (IPv4, URL, correo, fecha, UUID, trama de un sensor, entre otros) se verifican con cadenas que deben aceptarse y rechazarse.
  \item \textbf{Validación Lógica:} Se hace \emph{fuzzing}: se comparan los resultados del AFD mínimo con el módulo \texttt{re} de Python en 1500 cadenas aleatorias, que incluyen casos borde como la cadena vacía, saltos de línea y tabulaciones. Además, el código C generado se compila con \texttt{gcc} y se prueba directamente.
  \item \textbf{Consumo de Memoria:} Se compila el código generado de cuatro ejemplos junto con un programa de prueba y se lo ejecuta con Valgrind. En todos los casos el resumen informa \texttt{0 allocs, 0 frees, 0 bytes allocated}, sin errores de memoria. El programa de prueba no usa \texttt{stdio}, porque la biblioteca de C reserva en el \emph{heap} el buffer de la salida estándar en la primera escritura, y Valgrind contaría esa reserva aunque la función generada no haga ninguna.
  \item \textbf{Desempeño Comparativo:} Se traduce el AST de cada ejemplo a una expresión regular extendida de POSIX y se mide el tiempo medio de $10^6$ validaciones con la función generada y con \texttt{regexec()} de \texttt{regex.h}, sin contar la compilación de la expresión (\texttt{regcomp()}). Antes de medir se verifica que ambos clasifiquen igual 1000 cadenas por ejemplo (muestras y mutaciones aleatorias): no hubo ninguna discrepancia. La Tabla~\ref{tab:desempeno} muestra los resultados.
\end{enumerate}

\begin{table}[htbp]
\centering
\caption{Tiempo medio por cadena con $10^6$ validaciones (\texttt{gcc -O2}, x86-64). Los valores absolutos dependen de la máquina.}
\label{tab:desempeno}
\small
\setlength{\tabcolsep}{12pt}
\begin{tabular}{@{}lrrr@{}}
\toprule
\textbf{Ejemplo} & \textbf{NotFlex} & \textbf{\texttt{regex.h}} & \textbf{Aceleración} \\
\midrule
Identificador de C & 6,7 ns  & 76,5 ns  & 11,5$\times$ \\
IPv4               & 7,4 ns  & 62,1 ns  & 8,4$\times$ \\
Correo electrónico & 15,5 ns & 90,9 ns  & 5,9$\times$ \\
URL                & 35,8 ns & 134,6 ns & 3,8$\times$ \\
UUID               & 38,7 ns & 83,3 ns  & 2,2$\times$ \\
\bottomrule
\end{tabular}
\end{table}

La función generada fue entre 2,2 y 11,5 veces más rápida. La ventaja es mayor en los ejemplos de cadenas cortas (identificador, IPv4), lo que sugiere que allí pesa más el costo fijo de cada llamada a \texttt{regexec()}; en las cadenas largas domina el recorrido carácter por carácter de ambos.

% ============================================================================
\section{Comparación con Trabajos Relacionados}
% ============================================================================
Las librerías como \texttt{regex.h} o \texttt{PCRE2} trabajan en tiempo de ejecución. NotFlex, en cambio, sigue la idea de herramientas ya conocidas como \texttt{re2c} y \texttt{Ragel}, que hacen todo el trabajo durante la compilación (Tabla~\ref{tab:comparacion}).

El tiempo de ejecución de \texttt{regex.h} depende de la biblioteca que lo implemente. La de glibc, por ejemplo, no usa \emph{backtracking} salvo con retroreferencias: con patrones como \texttt{(a|aa)*b} o \texttt{(a+)+b}, que hacen explotar a los motores con \emph{backtracking}, medimos un tiempo lineal en cadenas de hasta 16\,000 caracteres. La diferencia de NotFlex es que el tiempo lineal está garantizado por construcción y no depende de la biblioteca de la plataforma.

\begin{table}[htbp]
\centering
\caption{Comparación de NotFlex frente a las alternativas tradicionales.}
\label{tab:comparacion}
\small
\renewcommand{\arraystretch}{1.15}
\begin{tabularx}{\textwidth}{@{}>{\raggedright\arraybackslash}p{3.1cm} *{3}{>{\centering\arraybackslash}X}@{}}
\toprule
\textbf{Característica} & \textbf{Validador a Mano} & \textbf{Librería (\texttt{regex.h})} & \textbf{NotFlex} \\
\midrule
\textbf{Dependencias externas} & No         & Sí                              & \textbf{No} \\
\textbf{Memoria dinámica}      & No         & Sí                              & \textbf{No ($0$ bytes)} \\
\textbf{Tiempo de ejecución}   & Variable   & Depende de la implementación & \textbf{Lineal $O(n)$} \\
\textbf{Mantenibilidad}        & Dificultosa & Alta                           & \textbf{Alta} \\

\bottomrule
\end{tabularx}
\end{table}

% ============================================================================
\section{Conclusiones}
% ============================================================================
\textbf{NotFlex} logró lo que buscábamos: una serie de pasos automatizados que se ejecutan de forma encadenada y van de la expresión regular al código C. Integramos las fases clásicas de un compilador (análisis léxico, sintáctico y semántico con tabla de símbolos, transformaciones de autómatas, minimización y generación de C). El resultado es una función en C que no usa memoria dinámica y siempre recorre la cadena en tiempo lineal, ideal para dispositivos con recursos limitados. Las pruebas lo confirman: Valgrind no registra ninguna reserva en el \emph{heap}, y la función generada coincide con \texttt{regex.h} en todas las cadenas probadas y es entre 2,2 y 11,5 veces más rápida.

% ============================================================================
\section*{Glosario}
% ============================================================================
\small
\begin{description}
  \item[\textbf{AFD}] Autómata Finito Determinista. En cada estado, para cada carácter, hay un solo camino posible.
  \item[\textbf{AFND}] Autómata Finito No Determinista. Puede tener varios caminos para un mismo carácter y transiciones épsilon.
  \item[\textbf{Análisis semántico}] Fase que verifica el significado de la entrada, por ejemplo que cada referencia tenga una definición y que no haya definiciones circulares.
  \item[\textbf{AST}] Árbol de Sintaxis Abstracta. Árbol que representa la estructura de la expresión regular después de analizarla.
  \item[\textbf{Backtracking}] Técnica que prueba un camino y, si falla, vuelve atrás a probar otro. Con ciertos patrones puede tardar muchísimo.
  \item[\textbf{C99}] Versión del estándar del lenguaje C publicada en 1999. Se optó por usar esta porque casi todos los compiladores de microcontroladores la soportan.
  \item[\textbf{Clase de equivalencia}] Grupo de caracteres que tienen las mismas transiciones en el autómata. Permite achicar la tabla.
  \item[\textbf{Cota}] Límite máximo. Si el tiempo de ejecución tiene cota, se asegura que nunca pasa de cierto valor.
  \item[\textbf{DSL}] \emph{Domain-Specific Language}. Lenguaje o herramienta hecha para resolver un solo tipo de problema, en este caso expresiones regulares.
  \item[\textbf{Embebido}] Sistema informático chico que está dentro de otro dispositivo y hace una tarea puntual, con poca memoria y poco procesamiento (por ejemplo, un sensor o un microondas).
  \item[\textbf{$\epsilon$-clausura}] Conjunto de estados a los que se llega desde un estado usando solamente transiciones épsilon (sin leer ningún carácter).
  \item[\textbf{Fuzzing}] Prueba que consiste en generar muchísimas entradas (incluso raras) y comparar los resultados contra una referencia para encontrar errores.
  \item[\textbf{Hopcroft}] Algoritmo que minimiza un AFD juntando los estados que no se pueden distinguir entre sí. El resultado es el AFD con menos estados posible.
  \item[\textbf{Host}] La PC donde se ejecuta el compilador, a diferencia del dispositivo donde después corre el código generado.
  \item[\textbf{IoT}] \emph{Internet of Things} (Internet de las Cosas). Objetos cotidianos con sensores que se conectan a una red.
  \item[\textbf{LL(1)}] Gramática que puede analizarse de izquierda a derecha decidiendo cada paso con un solo token de anticipación, lo que permite un parser descendente recursivo sin retroceso.
  \item[\textbf{MVP}] \emph{Minimum Viable Product}. Primera versión funcional con lo mínimo necesario.
  \item[\textbf{$O(n)$}] Tiempo lineal. El tiempo crece en proporción al largo $n$ de la cadena: si la cadena es el doble de larga, tarda el doble.
  \item[\textbf{Parser}] Parte del compilador que analiza el texto de la expresión regular y arma el AST.
  \item[\textbf{Pipeline}] Serie de pasos encadenados donde la salida de uno es la entrada del siguiente.
  \item[\textbf{Subconjuntos}] Algoritmo que convierte un AFND en un AFD. Cada estado del AFD representa un conjunto de estados del AFND.
  \item[\textbf{Tabla de símbolos}] Estructura que guarda cada nombre definido en la entrada junto con su información (aquí, su AST, su línea y la cantidad de usos).
  \item[\textbf{Thompson}] Algoritmo que convierte un AST en un AFND, armándolo por partes y uniéndolas con transiciones épsilon.
  \item[\textbf{Valgrind}] Herramienta que ejecuta un programa y detecta errores de memoria. La usamos para comprobar que no hay reservas dinámicas (\texttt{malloc}).
\end{description}
\normalsize

% ============================================================================
%  Referencias (formato splncs04)
% ============================================================================
\begin{thebibliography}{8}
\bibitem{aho2007}
Aho, A.V., Lam, M.S., Sethi, R., Ullman, J.D.: Compilers: Principles, Techniques, and Tools. Pearson Education, 2nd edn. (2007)

\bibitem{hopcroft2006}
Hopcroft, J.E., Motwani, R., Ullman, J.D.: Introduction to Automata Theory, Languages, and Computation. Addison-Wesley, 3rd edn. (2006)


\end{thebibliography}

\end{document}
